\documentclass[10pt,a4paper]{amsart}
\usepackage[margin=19mm]{geometry}
\usepackage{amsmath,amssymb,mathtools}
\usepackage[scr=rsfs,cal=euler]{mathalfa}
\usepackage{xfrac}
\usepackage[unicode,hidelinks,bookmarks=false]{hyperref}
\newcommand{\ie}{i.e.\ }
\newcommand{\cf}{cf.\ }
\newcommand{\resp}{respectively\ }
\newcommand{\Subsection}{\subsection}
\providecommand{\nobreakdash}{\nobreak\hbox{-}}
\setlength{\emergencystretch}{2em}
\multlinegap0pt
\usepackage{bm}
\usepackage{bbm}
\usepackage{dsfont}
\usepackage{xspace}
\usepackage{cancel}
\usepackage{mathtools}
\usepackage{amsthm}
\makeatletter
\@ifundefined{thm}{\newtheorem{thm}{Thm}}{}
\@ifundefined{lm}{\newtheorem{lm}[thm]{Lemma}}{}
\@ifundefined{prop}{\newtheorem{prop}[thm]{Proposition}}{}
\makeatother
\newcommand{\la}{\langle}
\newcommand{\pr}{\mathbb{P}}
\newcommand{\ra}{\rangle}
\title[On the Limit of the Tridiagonal Model for $\beta$-Dyson Brow]{On the Limit of the Tridiagonal Model for $\beta$-Dyson Brownian Motion}
\author{Alan Edelman, Sungwoo Jeong, Ron Nissim}
\date{Source: arXiv:2411.01633v3 (CC BY 4.0)}
\begin{document}
\maketitle

\noindent\emph{Source and licence.} On the Limit of the Tridiagonal Model for $\beta$-Dyson Brownian Motion. Version arXiv:2411.01633v3 (CC BY 4.0), distributed under the Creative Commons Attribution 4.0 International licence, as recorded in the arXiv licence metadata for this exact version and shown on the abstract page https://arxiv.org/abs/2411.01633v3. Full rights notice: licenses/arxiv-2411.01633-CC-BY-4.0.txt.\par\smallskip

\noindent\textit{Editorial scope.} Counted unit: the Prop whose source label is \texttt{vartheta = theta + epsilon} in the article (source0.tex lines 1032--1038, source label \texttt{vartheta = theta + epsilon}), delivered together with the article's own complete original proof of it (source0.tex lines 1040--1148, one \texttt{proof} environment of 109 lines). No mathematics is written, repaired or re-derived: statement and proof are the article's own text line for line. Required context retained uncounted: Three term reccurence (lines 473--476); Concentration for specific functionals of theta (lines 983--1001). Every definition, display and label of those lines is retained unchanged. Omitted from this package and not redistributed: 1--472, 477--982, 1002--1031, 1039--1039, 1149--1642. Nothing inside a retained excerpt is omitted or altered: every retained body is delivered exactly as the article's lines are written, so no omission receipt of any kind accompanies this package. Citations: the retained excerpts carry no citation command at all, so no bibliography entry of the article is transcribed and no reference is redistributed. Reference adaptation: none was needed and none was made; every reference inside the delivered unit resolves inside it, so no unresolved reference or citation is produced. Printed slips kept literal: no printed word, symbol, hypothesis or number of the counted statement or of its proof was altered. Layout and class: the article's own class is \texttt{amsart} with packages including geometry, amsmath, amsfonts, amssymb, natbib, graphicx, enumerate, mathtools, subcaption, xcolor, comment, tikz, bbm, mathrsfs; the wrapper is the lane's pinned \texttt{amsart} editorial wrapper at 10pt on A4, which restates the theorem-like environments the retained text needs and adds the article's own macro definitions that the retained text needs (\texttt{\textbackslash la}, \texttt{\textbackslash pr}, \texttt{\textbackslash ra}), with extra wrapper packages (bm, bbm, dsfont, xspace, cancel, mathtools). The delivered numbering is the unit's own. Assets: no retained span contains a figure environment, a graphics-inclusion command, a PDF-page inclusion, a table or a TikZ picture; no image file is copied into this package, the package declares no asset dependency and no image file is redistributed with it. Licence: version arXiv:2411.01633v3 of this article is distributed under the Creative Commons Attribution 4.0 International licence, as recorded in the arXiv licence metadata for this exact version and shown on the abstract page https://arxiv.org/abs/2411.01633v3. A full rights notice is delivered at licenses/arxiv-2411.01633-CC-BY-4.0.txt. Characters: every retained excerpt is pure ASCII, so the delivered engine, XeLaTeX with its default Latin Modern text font, reports no missing character.
\begin{equation}
    P_{k+1}(x) = x P_k(x)-P_{k-1}(x).
    \label{Three term reccurence}
\end{equation}

\begin{lm}\label{Concentration for specific functionals of theta}
    There is a $C=C(T,k)>0$ such that on a set $A_k$ with $\pr(A_k)>1-C'(T,k)n^{-2}$, the following holds. 
    \begin{enumerate}
        \item For any deterministic unit vector $v \in V_{k}:=\mathrm{Span}(\{e_j\}_{j=1}^k)$ and $x(t) \in \{\theta_j(t)\}_{j\leq k} \cup\{M(t)^{j_1}e_{j_2+2}\}_{0 \leq j_1 \leq k, 0 \leq j_2 \leq k-1}$, we have
    \begin{equation}
        \sup_{t \in [0,T]} |\la x(t),v \ra| \leq \frac{C\log(n)}{\sqrt{n}}.
        \label{Single Component bound}
    \end{equation}
    \item For $j, \ell =0,1,\dots,k$,
    \begin{equation}
        \sup_{t \in [0,T]} |\theta_j(t)^T\theta_{\ell}(t)-\delta_{j,\ell}| \leq \frac{C\log(n)}{\sqrt{n}}.
        \label{inner product concentration}
    \end{equation}
    \item 
    \begin{equation}
        \sup_{t \in [0,T]} \max(\|M(t)\|_{op},\|\tilde{M}(t)\|_{op}) \leq C\log(n).
    \end{equation}
    \end{enumerate}
\end{lm}

\begin{prop}
   On the event $A_k$ defined in Lemma \ref{Concentration for specific functionals of theta}, for $n$ sufficiently large, we have,
\begin{equation}\label{vartheta = theta + epsilon}
        \vartheta_k(t) = (c_k(t)+\mu_k(t))\theta_k(t)+\epsilon_k(t),
    \end{equation}
    with $c_k(t):=\frac{1}{\|\theta_{k-1}(t)\|}$, $\epsilon_k(t) \in \Theta_{k-1}(t)+ \mathcal{E}_k(t)$, $\sup_{t \in [0,T]}\|\epsilon_k(t)\| \leq C_{k,T} \frac{\log(n)}{\sqrt{n}}$, and $\sup_{t \in [0,T]} |\mu_k(t)| \leq C_{k,T}\frac{\log(n)^{2(k-1)}}{n}$.
\end{prop}

\begin{proof}
The proof is by the strong induction on $k$. The base case $k=0$ is trivial. For the inductive step, we start by proving that on the event $A_k$ defined in Lemma \ref{Concentration for specific functionals of theta}, for $n$ sufficiently large, we have,
\begin{equation}\label{Householder Product}
        \begin{split}
            &Q_0(t)Q_1(t)\cdots Q_{k-1}(t)\\
            &=I +\sum_{i < k, 2\leq j \leq k} \left(c_{i,j}(t) \vartheta_i(t)e_j^T+d_{i,j}(t) e_j\vartheta_i(t)^T\right)+\sum_{j,j' < k} f_{j,j'}(t) \vartheta_j(t)\vartheta_{j'}(t)^T,
        \end{split}
    \end{equation}
for some processes $c_{i,j}(t), d_{i,j}(t), f_{j,j'(t)}$, with $\max_{i,j,j'} \sup_{t \in [0,T]} (|c_{i,j}(t)|+ |d_{i,j}(t)|+|f_{j,j'(t)}|) \leq C_{k,T}$ for some constant $C_{k,T}>0$ depending only on $k$ and $T$. 


Observe that by the inductive hypothesis and Lemma \ref{Concentration for specific functionals of theta}, it is clear that $|\vartheta_j(t)^T \vartheta_{j'}(t)-\delta_{j,j'}| \leq C_k \frac{\log(n)}{\sqrt{n}}$ and $|\vartheta_j(t)^T e_{j'}|\leq C_k \frac{\log(n)}{\sqrt{n}}$ for $j, j' \leq k$. As a result, we obtain
\begin{equation*}
   \left| \frac{2}{\|\vartheta_j(t)-\|\vartheta_j(t)\|e_{k+2}\|^2}\right| \leq 2+C_k\frac{\log(n)}{\sqrt{n}}.
\end{equation*}
Expanding the product $Q_0(t)Q_1(t)\cdots Q_{k}(t)$, an expansion of the form in the right hand side of \eqref{Householder Product} is clear. Moreover, the $c_{i,j}(t),d_{i,j}(t), f_{j,j'}(t)$ are all polynomials of $\frac{2}{||\vartheta_j(t)-||\vartheta_j(t)||e_{k+2}||^2}$, $e_a^T \vartheta_b(t)$ $\vartheta_a(t)^T \vartheta_b(t)$, and $\delta_a(t)$, for $a,b \leq k$, with number of terms and coefficients only depending on $k$. This completes the proof of \eqref{Householder Product}. The proof of the inductive step will involve repeated applications of \eqref{Householder Product} in tandem with the induction hypothesis and Lemma \ref{Concentration for specific functionals of theta}.

For the inductive step, we seek to prove  \eqref{vartheta = theta + epsilon} in the $k+1$ case. First, from the construction of the Householder matrix, $Q_k(t)e_{k+2}=\frac{\vartheta_k(t)}{||\vartheta_k(t)||}$. Once again recalling that $|\vartheta_j(t)^T \vartheta_{j'}(t)-\delta_{j,j'}| \leq C_{k,T} \frac{\log(n)}{\sqrt{n}}$ and $|\vartheta_j(t)^T e_{j'}|\leq C_{k,T} \frac{\log(n)}{\sqrt{n}}$ for $j, j' \leq k$, we have 

\begin{equation*}
    \begin{split}
        ~&Q_0(t)\cdots Q_{k-1}(t)Q_k(t)e_{k+2}\\
        &= \bigg(I+\sum_{i <k, j \leq k} (c_{i,j}(t) \vartheta_i(t)e_j^T+d_{i,j}(t) e_j\vartheta_i(t)^T)+\sum_{j,j' < k } f_{j,j'}(t) \vartheta_j(t)\vartheta_{j'}(t)^T\bigg) \frac{\vartheta_k(t)}{||\vartheta_k(t)||}\\
            & = \frac{\vartheta_k(t)}{||\vartheta_k(t)||}+\varepsilon_k(t),
    \end{split}
\end{equation*}
where $||\varepsilon_k(t)||\leq C_k\frac{\log(n)}{\sqrt{n}}$, and $\varepsilon_k(t) \in \mathrm{Span}\{\vartheta_j(t)\}_{j\leq k-1}+\mathrm{Span}\{e_j(t)\}_{2\leq j\leq k+1} \subset \Theta_{k-1}(t)+\mathcal{E}_{k}(t)$. Applying the induction hypothesis to the last equation display, we can write
\begin{equation*}
    Q_0(t)\cdots Q_{k-1}(t)Q_k(t)e_{k+2}  =\left(c_k(t)+\mu_k(t)\right)\frac{\theta_k(t)}{||\vartheta_{k}(t)||}+\varepsilon_k'(t),
\end{equation*}
where $||\varepsilon_k'(t)||\leq C_{k,T}\frac{\log(n)^k}{\sqrt{n}}$, and $\varepsilon_k'(t) \in \Theta_{k-1}(t)+\mathcal{E}_{k}(t)$. 

Using the three term recurrence for the Chebyshev polynomials \eqref{Three term reccurence} and the norm bound of Lemma \ref{Concentration for specific functionals of theta}, we have,
\begin{equation*}
    \begin{split}
        &\frac{1}{\sqrt{n}}\Tilde{M}(t)Q_0(t)\cdots Q_{k-1}(t)Q_k(t)e_{k+2} \\
        &= \frac{c_k(t)+\mu_k(t)}{||\vartheta_{k}(t)||}(\theta_{k+1}(t)+\theta_{k-1}(t))+ \varepsilon_{k+1}(t)\\
        &=\frac{c_k(t)+\mu_k(t)}{||\vartheta_{k}(t)||}(\theta_{k+1}(t)+\vartheta_{k-1}(t))+ \varepsilon_{k+1}'(t),
    \end{split}
\end{equation*}
where $||\varepsilon_{k+1}(t)||,||\varepsilon_{k+1}'(t)||\leq C_{k,T}\frac{\log(n)^{k+1}}{\sqrt{n}}$, and $\varepsilon_{k+1}(t), \varepsilon_{k+1}'(t) \in \Theta_{k}(t)+\Tilde{M}(t)\mathcal{E}_{k}(t) \subseteq \Theta_{k}(t)+\mathcal{E}_{k+1}(t)$. Then once again using the induction hypothesis, Lemma \ref{Concentration for specific functionals of theta}, and equation \eqref{Householder Product}, we get
\begin{equation*}
    \begin{split}
        &\frac{1}{\sqrt{n}}Q_{k-2}(t)\cdots Q_{0}(t)\Tilde{M}(t)Q_0(t)\cdots Q_{k-1}(t)Q_k(t)e_{k+2}\\
        &\hspace{2cm}= \frac{c_k(t)+\mu_k(t)}{||\vartheta_{k}(t)||}(\theta_{k+1}(t)+\vartheta_{k-1}(t))+ \varepsilon_{k+1}''(t),
    \end{split}
\end{equation*}
where $||\varepsilon_{k+1}''(t)||\leq C_{k,T}\frac{\log(n)^{k+1}}{\sqrt{n}}$, and $\varepsilon_{k+1}''(t) \in \Theta_{k}(t)+\mathcal{E}_{k}(t)+\Tilde{M}(t)\mathcal{E}_{k}(t) \subseteq \Theta_{k}(t)+\mathcal{E}_{k+1}(t)$. Recall that $Q_{k-1}(t)\vartheta_{k-1}(t) =||\vartheta_{k-1}(t)||e_{k+1}$, so applying $Q_{k-1}(t)$ on the left we have
\begin{equation*}
    \begin{split}
        &\frac{1}{\sqrt{n}}Q_{k-1}(t)\cdots Q_{0}(t)\Tilde{M}(t)Q_0(t) \cdots Q_{k-1}(t)Q_k(t)e_{k+2} \\
        &\hspace{2cm}= \frac{c_k(t)+\mu_k(t)}{||\vartheta_{k}(t)||}(\theta_{k+1}(t)+||\vartheta_{k-1}(t)||e_{k+1})+ \varepsilon_{k+1}'''(t),
    \end{split}
\end{equation*}
where $||\varepsilon_{k+1}'''(t)||\leq C_{k,T}\frac{\log(n)^{k+1}}{\sqrt{n}}$, and $\varepsilon_{k+1}'''(t) \in \Theta_{k}(t)+\mathcal{E}_{k+1}(t)$. Finally, applying $Q_k(t)$ on the left we obtain
\begin{equation*}
    \begin{split}
        &\frac{1}{\sqrt{n}}Q_k(t)Q_{k-1}(t)\cdots Q_{0}(t)\Tilde{M}(t)Q_0(t)\cdots Q_{k-1}(t)Q_k(t)e_{k+2} \\
        &\hspace{2cm} = \frac{c_k(t)+\mu_k(t)}{||\vartheta_{k}(t)||}(\theta_{k+1}(t)+||\vartheta_{k-1}(t)||e_{k+1})+ \varepsilon_{k+1}^{(4)}(t),
    \end{split}
\end{equation*}
with $||\varepsilon_{k+1}^{(4)}(t)||\leq C_{k,T}\frac{\log(n)}{\sqrt{n}}$, and $\varepsilon_{k+1}^{(4)}(t) \in \Theta_{k}(t)+\mathcal{E}_{k+1}(t)$. 

For the next step we estimate the error from $M(t)\approx \Tilde{M}(t)$. Using \eqref{Householder Product} and the fact that the first $k$ entries of $\vartheta_k(t)$ are zero by definition, we first see that
\begin{equation}\label{eq:Q's-fix-basis}
    Q_k(t)Q_{k-1}(t)\cdots Q_{0}(t)e_1=e_1.
\end{equation}
So plugging in the definition $\tilde{M}(t)=(I-e_1e_1^T)M(t)(I-e_1e_1^T)$ and applying \eqref{eq:Q's-fix-basis},
\begin{equation*}
    \begin{split}
        ~&\frac{1}{\sqrt{n}}Q_k(t)Q_{k-1}(t)\cdots Q_{0}(t)(M(t)-\tilde{M}(t))Q_0(t)\cdots Q_{k-1}(t)Q_k(t)\\
        &=\frac{1}{\sqrt{n}}Q_{k}(t)\cdots Q_{0}(t)\left(e_1 e_1^T M(t)+M(t) e_1e_1^T-e_1e_1^TM(t)e_1e_1^T\right)\\
        &Q_0(t)\cdots Q_{k-1}(t)Q_{k-1}(t)\\
        &=\frac{1}{\sqrt{n}}(e_1e_1^T M(t)Q_0(t)\cdots Q_{k-1}(t)Q_{k-1}(t)\\
        &+(e_1e_1^T M(t)Q_0(t)\cdots Q_{k-1}(t)Q_{k-1}(t))^T+ e_1e_1^T M(t)e_1 e_1^T) \\
        &=\frac{1}{\sqrt{n}}(e_1e_1^T M(t)+M(t)e_1e_1^T+e_1e_1^T M(t)e_1 e_1^T)+E_k(t)\\ 
        &= \frac{1}{\sqrt{n}}(M(t)-\tilde{M}(t))+E_k(t),
    \end{split}
\end{equation*}
where $\|E_k(t)\|\leq C_k\frac{\log(n)}{\sqrt{n}}$, and $\mathrm{Range}(E_k(t)) \subset \Theta_{k-1}(t)+\mathcal{E}_{k+1}(t)$. These properties of the error term in the above equation display, $E_k(t)$, follow from equation \eqref{Householder Product} and Lemma \ref{Concentration for specific functionals of theta} once again. In particular, we have
\begin{equation*}
    \|\frac{1}{\sqrt{n}}Q_{k}(t)\cdots Q_{0}(t)(M(t)-\Tilde{M}(t))Q_0(t)\cdots Q_{k-1}(t)Q_k(t)e_{k+2}\| \leq C_{k,T}\frac{\log(n)}{\sqrt{n}},
\end{equation*}
with $\frac{1}{\sqrt{n}}Q_{k}(t)\cdots Q_{0}(t)(M(t)-\Tilde{M}(t))Q_0(t)\cdots Q_{k-1}(t)Q_k(t)e_{k+2} \in \Theta_{k}(t)+\mathcal{E}_{k+1}(t)$. 

Combining everything so far, we have
\begin{equation*}
    \begin{split}
        \frac{1}{\sqrt{n}}M_{k+1}(t)e_{k+2}=\frac{c_k(t)+\mu_k(t)}{||\vartheta_{k}(t)||}\theta_{k+1}(t)+ \varepsilon_{k+1}^{(5)}(t),
    \end{split}
\end{equation*}
with $\|\varepsilon_{k+1}^{(5)}(t)\|\leq C_{k,T}\frac{\log(n)^{k+1}}{\sqrt{n}}$, and $\varepsilon_k^{(5)}(t) \in \Theta_{k}(t)+\mathcal{E}_{k}(t)+\Tilde{M}(t)\mathcal{E}_{k}(t)$. However, by the induction hypothesis we have
\begin{equation}\label{Norm of theta computation}
    \begin{split}
        ||\vartheta_k(t)||^2&=(c_k(t)+\mu_k(t))^2||\theta_k(t)||^2+2(c_k(t)+\mu_k(t))\epsilon_k(t)^T\theta_k(t)+||\epsilon_k(t)||^2\\
        &=c_k(t)^2||\theta_k(t)||^2+O(\log(n)^{2k}/n),
    \end{split}
\end{equation}
and thus similarly $\frac{1}{\|\vartheta_k(t)\|}=\frac{1}{c_k(t)\|\theta_k(t)\|}+O(\log(n)^{2k}/n)$. Finally we deduce that 
\begin{equation*}
    \begin{split}
        \vartheta_{k+1}(t)&=\frac{1}{\sqrt{n}}M_{k+1}(t)e_{k+2}-\|\vartheta_{k-1}(t)\|e_{k+1}-\frac{a_{k+1}(t)}{\sqrt{n}}\\
        &=\bigg(\frac{1}{\|\theta_{k}(t)\|}+O(\log(n)^{2k}/n)\bigg)(\theta_{k+1}(t)+\|\vartheta_{k-1}(t)\|e_{k+1}+ \epsilon_{k+1}'(t))\\
        &\hspace{7cm}-\|\vartheta_{k-1}(t)\|e_{k+1}-\frac{a_{k+1}(t)}{\sqrt{n}}\\
        &=c_{k+1}(t)\theta_{k+1}(t)+\mu_{k+1}(t)\theta_{k+1}(t)+\epsilon_{k+1}(t).
    \end{split}
\end{equation*}
With $c_{k+1}(t),\mu_{k+1}(t),\epsilon_{k+1}(t)$ all satisfying the desired properties. Here we used the fact that $|a_{k+1}(t)| \leq C_{k,T}$ which is easy to prove by the induction hypothesis and Lemma \ref{Concentration for specific functionals of theta}. This completes the proof. 
\end{proof}

\end{document}
